% \iffalse meta-comment
%
% Copyright (C) 2019-2025 by Jan Hajer
% -----------------------------------
%
% This file may be distributed and/or modified under the
% conditions of the LaTeX Project Public License, either version 1.3c
% of this license or (at your option) any later version.
% The latest version of this license is in:
%
% http://www.latex-project.org/lppl.txt
%
% and version 1.3c or later is part of all distributions of LaTeX
% version 2005/12/01 or later.
%
% \fi
%
% \iffalse

%<package>\NeedsTeXFormat{LaTeX2e}[2021-11-15]
%<package>\ProvidesPackage{hep-cv-private}[2026/08/01 v1.4 Curricula Vitarum in High Energy Physics]
%
%<*documentation>
\RequirePackage[l2tabu, orthodox]{nag}
\documentclass{ltxdoc}

\renewcommand*\theCodelineNo{\rmfamily\tstyle\footnotesize\arabic{CodelineNo}}
\AtBeginEnvironment{macrocode}{\renewcommand*{\ttdefault}{clmt}}
\renewcommand*{\MacroFont}{\codestyle}
\AtBeginDocument{\DeleteShortVerb{\|}}
\AtBeginDocument{\MakeShortVerb{\"}}
\EnableCrossrefs
\CodelineIndex
\RecordChanges

\usepackage{hologo}

\usepackage[parskip,oldstyle,font=10pt]{hep-paper}
\bibliography{bibliography}
\acronym{CV}{curriculum vitae}[curricula vitarum]
\acronym{TOC}{table of contents}
%</documentation>

%<*driver>
\expandafter\newif\csname ifshort\endcsname
\shortfalse
\begin{document}
\DocInput{hep-cv-private-implementation.dtx}
\end{document}
%</driver>
%
% \fi
%
% \CheckSum{251}
%
% \CharacterTable
%  {Upper-case    \A\B\C\D\E\F\G\H\I\J\K\L\M\N\O\P\Q\R\S\T\U\V\W\X\Y\Z
%   Lower-case    \a\b\c\d\e\f\g\h\i\j\k\l\m\n\o\p\q\r\s\t\u\v\w\x\y\z
%   Digits        \0\1\2\3\4\5\6\7\8\9
%   Exclamation   \!     Double quote  \"     Hash (number) \#
%   Dollar        \$     Percent       \%     Ampersand     \&
%   Acute accent  \'     Left paren    \(     Right paren   \)
%   Asterisk      \*     Plus          \+     Comma         \,
%   Minus         \-     Point         \.     Solidus       \/
%   Colon         \:     Semicolon     \;     Less than     \<
%   Equals        \=     Greater than  \>     Question mark \?
%   Commercial at \@     Left bracket  \[     Backslash     \\
%   Right bracket \]     Circumflex    \^     Underscore    \_
%   Grave accent  \`     Left brace    \{     Vertical bar  \|
%   Right brace   \}     Tilde         \~}
%
% \changes{v1.0}{2019/01/01}{Created the private CV bibliography filters and environments for work, visits, education, supervision, funding, talks, publications, software, theses, and service entries.}
% \changes{v1.3}{2024/11/01}{Split the private CV helpers into a dedicated documented package with keyword filters, CV bibliography environments, and Indico e-print formatting.}
% \changes{v1.4}{2026/08/01}{Raised the LaTeX format requirement, bumped release metadata, replaced the stub documentation with package documentation, added README content, and added README extraction to the installer.}
%
% \ifshort
%<*documentation>
% \fi
%
\GetFileInfo{hep-cv-private.sty}

\title{The \software{hep-cv-private} package\thanks{This document corresponds to \software{hep-cv-private}~\fileversion.}}

\author{Jan Hajer \email{jan.hajer@tecnico.ulisboa.pt}}

\date{\filedate}

% \ifshort
\begin{document}
% \fi

\maketitle*

\begin{abstract}
The \software{hep-cv-private} package adds private CV bibliography filters and environments on top of the HEP CV bibliography setup.
It groups work, visits, education, supervision, funding, talks, publications, software, theses, and service entries by keyword, provides specialised CV bibliography environments for those groups, and formats Indico event and contribution links as eprint data.
\end{abstract}

\tableofcontents\clearpage

\section{Introduction}

Load the package with "\usepackage{hep-cv-private}" after the CV bibliography setup.

\DescribeMacro{work}
\DescribeMacro{visit}
\DescribeMacro{education}
\DescribeMacro{lecture}
\DescribeMacro{tutorial}
\DescribeMacro{master-supervision}
\DescribeMacro{PhD}
\DescribeMacro{master}
\DescribeMacro{funding-individual}
\DescribeMacro{funding-project}
\DescribeMacro{funding-participation}
\DescribeMacro{organization}
The package defines bibliography filters for CV material using keywords in the bibliography database.
The filters cover work, visits, education, lectures, tutorials, master supervision, PhD and master entries, individual, project, and participation funding, and organisation or service entries.

\DescribeMacro{seminar}
\DescribeMacro{conference}
\DescribeMacro{contributions}
\DescribeMacro{student}
Talk-related filters separate invited seminars, invited conferences, contributed conference, workshop, or school presentations, and student presentations.

\DescribeMacro{preprint}
\DescribeMacro{journal}
\DescribeMacro{proceedings}
\DescribeMacro{editor}
\DescribeMacro{community}
\DescribeMacro{software}
\DescribeMacro{thesis}
\DescribeMacro{main}
Publication filters group preprints, journal articles, proceedings, editor entries, community entries, software, theses, and selected main-recent publications.

\DescribeEnv{work}
\DescribeEnv{visit}
\DescribeEnv{education}
\DescribeEnv{student}
\DescribeEnv{course}
\DescribeEnv{funding}
\DescribeEnv{talk}
\DescribeEnv{publication}
\DescribeEnv{main}
The corresponding bibliography environments customise the CV left-column layout and field formatting for positions, visits, education, students, courses, funding, talks, publications, and selected main publications.
They move titles, dates, event information, funding titles, annotations, or presentation styles to the positions expected in a private CV.

\DescribeMacro{CERN}
\DescribeMacro{INFN}
\DescribeMacro{IJS}
\DescribeMacro{DESY}
\DescribeMacro{FNAL}
\DescribeMacro{Edinburgh}
\DescribeMacro{CNRS}
\DescribeMacro{IIHE}
\DescribeMacro{IRMP}
\DescribeMacro{MITP}
\DescribeMacro{LIP}
The package also defines Indico eprint formats for common event servers.
Entries can use these eprint types to link both the event and a contribution number.

% \ifshort
\printbibliography

\end{document}
%
%</documentation>
% \fi
%
% \StopEventually{
% \printbibliography
% \PrintChanges
% }
%
% \appendix
%
% \section{Implementation}
%
%<*package>
%
% \subsection{Bibliography filters}
%
% For CV
%    \begin{macrocode}
\defbibfilter{work}{keyword=work}
\defbibfilter{visit}{keyword=visit}
\defbibfilter{education}{keyword=education}
\defbibfilter{lecture}{keyword=lecture}
\defbibfilter{tutorial}{keyword=tutorial}
\defbibfilter{master-supervision}{keyword=master and keyword=supervision}
\defbibfilter{PhD}{keyword=PhD and not keyword=supervision}
\defbibfilter{master}{keyword=master and not keyword=supervision}
\defbibfilter{funding-individual}{keyword=funding and keyword=individual}
\defbibfilter{funding-project}{keyword=funding and keyword=project}
\defbibfilter{funding-participation}{keyword=funding and keyword=participation}
\defbibfilter{organization}{keyword=organization}
%    \end{macrocode}
% For talks
%    \begin{macrocode}
\defbibfilter{seminar}{keyword=invited
  and ( keyword=seminar or keyword=colloquium )}
\defbibfilter{conference}{keyword=invited and not keyword=organization
  and ( keyword=conference or keyword=workshop or keyword=school )}
\defbibfilter{contributions}{not keyword=invited and not keyword=organization
  and ( keyword=conference or keyword=workshop or keyword=school )}
\defbibfilter{student}{not keyword=invited and not keyword=organization and
  ( keyword=conference or keyword=workshop or keyword=school ) and keyword=student}
%    \end{macrocode}
% For publications
%    \begin{macrocode}
\defbibfilter{preprint}{keyword=preprint}
\defbibfilter{journal}{keyword=article}
\defbibfilter{proceedings}{keyword=proceedings}
\defbibfilter{editor}{keyword=editor}
\defbibfilter{community}{keyword=community and not keyword=editor}
\defbibfilter{software}{keyword=software}
\defbibfilter{thesis}{keyword=thesis}
\defbibfilter{main}{keyword=mainrecent}
%    \end{macrocode}
%
% \subsection{Macros}
%
% \begin{macro}{presentationstyle}
% Presentation mode.
%    \begin{macrocode}
\newbibmacro{presentationstyle}{%
  \ifkeyword{lecture}{Blackboard lecture}{%
  \ifkeyword{plenary}{Plenary talk}{%
  \ifkeyword{parallel}{Talk}{%
  \ifkeyword{talk}{Talk}{%
  \ifkeyword{poster}{Poster}{%
  \ifkeyword{seminar}{Seminar}{%
  \ifkeyword{colloquium}{Colloquium}{%
  \ifkeyword{convener}{Session convener}{}}}}}}}}%
}
%    \end{macrocode}
% \end{macro}
%
% \subsection{Bibliography environments}
%
% \begin{environment}{work}
% Work bibliography environment.
%    \begin{macrocode}
\defbibenvironment{work}{%
  \toggletrue{blx@skiplab}%
  \DeclareFieldFormat*{title}{\textbf{\sentence##1}}%
  \adjustOrgLocDat%
  \list{\printcvdate}{\setlengths}%
  \renewcommand*{\makelabel}[1]{\hss##1}%
}{\endlist}{\item}
%    \end{macrocode}
%\end{environment}
%
% \begin{environment}{visit}
% Visit bibliography environment.
%    \begin{macrocode}
\defbibenvironment{visit}{%
  \toggletrue{blx@skiplab}%
  \DeclareListFormat{organization}{\textbf{\sentence##1}}%
  \adjustOrgLocDat%
  \list{\printcvdate}{\setlengths}%
  \renewcommand*{\makelabel}[1]{\hss##1}%
}{\endlist}{\item}
%    \end{macrocode}
%\end{environment}
%
% \begin{environment}{education}
% Education bibliography environment.
%    \begin{macrocode}
\defbibenvironment{education}{%
  \toggletrue{blx@skiplab}%
  \DeclareFieldFormat*{title}{\textbf{\sentence##1}}%
  \adjustOrgLocDat%
  \list{\printcvdate}{\setlengths}%
  \renewcommand*{\makelabel}[1]{\hss##1}%
}{\endlist}{\item}
%    \end{macrocode}
%\end{environment}
%
% \begin{environment}{student}
% Student bibliography environment.
%    \begin{macrocode}
\defbibenvironment{student}{%
  \togglefalse{abx@bool@giveninits}%
  \toggletrue{blx@skiplab}%
  \DeclareNameAlias{author}{given-family-bold}%
  \adjustOrgLocDat%
  \letbibmacro{finentry-old}{finentry}%
  \renewbibmacro{finentry}{\usebibmacro{annotation}\usebibmacro{finentry-old}}%
  \list{\printcvdate}{\setlengths}%
  \renewcommand*{\makelabel}[1]{\hss##1}%
}{\endlist}{\item}
%    \end{macrocode}
%\end{environment}
%
% \begin{environment}{course}
% Course bibliography environment.
%    \begin{macrocode}
\defbibenvironment{course}{%
  \togglefalse{abx@bool@giveninits}%
  \toggletrue{blx@skiplab}%
  \adjustOrgLocDat%
  \DeclareFieldFormat*{title}{\textbf{\sentence##1}\newunit}%
  \list{\printissuedate}{\setlengths}%
  \renewcommand*{\makelabel}[1]{\hss##1}%
}{\endlist}{\item}
%    \end{macrocode}
%\end{environment}
%
% \begin{environment}{funding}
% Funding bibliography environment.
%    \begin{macrocode}
\defbibenvironment{funding}{%
  \letbibmacro{begentry-old}{begentry}%
  \renewbibmacro*{begentry}{%
    \usebibmacro{begentry-old}\usebibmacro{booktitle}\newunit%
  }%
  \DeclareFieldFormat*{booktitle}{\textbf{\sentence##1}}%
  \renewbibmacro*{author}{}%
  \renewbibmacro*{date}{}%
  \DeclareFieldFormat*{title}{\emph{Proposal title: ##1.}}%
  \list{\printindex\printcvdate}{\setlengths}%
  \renewcommand*{\makelabel}[1]{\hss##1}%
}{\endlist}{\item}
%    \end{macrocode}
%\end{environment}
%
% \begin{environment}{talk}
% Talk bibliography environment.
%    \begin{macrocode}
\defbibenvironment{talk}{%
  \toggletrue{blx@skiplab}%
  \letbibmacro{begentry-old}{begentry}
  \renewbibmacro*{begentry}{%
    \usebibmacro{begentry-old}\usebibmacro{booktitle}\newunit%
  }%
  \adjustOrgLocDat%
  \renewbibmacro{author}{}
  \DeclareFieldFormat{title}{%
    \emph{\usebibmacro{presentationstyle}: ##1}.
  }%
  \DeclareFieldFormat{booktitle}{\textbf{\sentence##1}}%
  \list{\printcvdate}{\setlengths}%
  \renewcommand*{\makelabel}[1]{\hss##1}%
}{\endlist}{\item}
%    \end{macrocode}
%\end{environment}
%
% \begin{environment}{publication}
% Publication bibliography environment.
%    \begin{macrocode}
\defbibenvironment{publication}{%
  \DeclareFieldFormat*{title}{\textbf{\sentence##1}\newunit}%
  \letbibmacro{begentry-old}{begentry}%
  \renewbibmacro*{begentry}{%
    \usebibmacro{begentry-old}\usebibmacro{titlefirst}\newunit%
  }%
  \renewbibmacro*{title}{}%
  \renewbibmacro*{in:}{%
    \setunit{\addcomma\space}%
    \iffieldundef{journaltitle}{}{%
      \ifkeyword{accepted}{\printtext{Accepted in}}{%
        \ifkeyword{submitted}{\printtext{Submitted to}}{%
          \printtext{\bibstring{in}}%
        }%
      }%
    }\intitlepunct%
  }%
  \renewbibmacro*{date}{}%
  \renewbibmacro{issue+date}{}%
  \list{\printindex\printcvdate}{\setlengths}%
  \renewcommand*{\makelabel}[1]{\hss##1}%
}{\endlist}{\item}
%    \end{macrocode}
%\end{environment}
%
% \begin{environment}{main}
% Key publication bibliography environment.
%    \begin{macrocode}
\defbibenvironment{main}{%
  \toggletrue{blx@skiplab}%
  \DeclareFieldFormat*{title}{\textbf{\sentence##1}\newunit}%
  \letbibmacro{begentry-old}{begentry}%
  \renewbibmacro*{begentry}{%
    \usebibmacro{begentry-old}\usebibmacro{titlefirst}\newunit%
  }%
  \renewbibmacro*{title}{}%
  \renewbibmacro*{issue+date}{}%
  \DeclareFieldFormat{annotation}{\newline ##1}%
  \letbibmacro{finentry-old}{finentry}%
  \renewbibmacro{finentry}{\usebibmacro{annotation}\usebibmacro{finentry-old}}%
  \list{\printcvdate}{\setlengths}%
  \renewcommand*{\makelabel}[1]{\hss##1}%
}{\endlist}{\item}
%    \end{macrocode}
%\end{environment}
%
% \subsection{Indico as eprint}
%
% \begin{macro}{\hep@indico}
% Define new indico command
%    \begin{macrocode}
\NewDocumentCommand{\hep@indico}{mm}{
  \DeclareFieldFormat{eprint:#1}{%
    \newcommand*{\hep@path}{event/\thefield{eprintclass}}%
    \mkbibacro{URL}\addcolon\space%
    \hep@bib@online{https://#2/\hep@path}{#2/\hep@path}%
    \hep@bib@online{https://#2/\hep@path/contributions/##1}{/contributions/##1}%
  }%
}
%    \end{macrocode}
% Use new indico command.
%    \begin{macrocode}
\hep@indico{CERN}{indico.cern.ch}
\DeclareFieldAlias{eprint:cern}{eprint:CERN}
\hep@indico{INFN}{agenda.infn.it}
\DeclareFieldAlias{eprint:infn}{eprint:INFN}
\hep@indico{IJS}{indico.ijs.si}
\DeclareFieldAlias{eprint:ijs}{eprint:IJS}
\hep@indico{DESY}{indico.desy.de}
\DeclareFieldAlias{eprint:desy}{eprint:DESY}
\hep@indico{FNAL}{indico.fnal.gov}
\DeclareFieldAlias{eprint:fnal}{eprint:FNAL}
\hep@indico{Edinburgh}{indico.ph.ed.ac.uk}
\DeclareFieldAlias{eprint:edinburgh}{eprint:Edinburgh}
\hep@indico{CNRS}{indico.in2p3.fr}
\DeclareFieldAlias{eprint:cnrs}{eprint:CNRS}
\hep@indico{IIHE}{indico.iihe.ac.be}
\DeclareFieldAlias{eprint:iihe}{eprint:IIHE}
\hep@indico{IRMP}{agenda.irmp.ucl.ac.be}
\DeclareFieldAlias{eprint:irmp}{eprint:IRMP}
\hep@indico{MITP}{indico.mitp.uni-mainz.de}
\DeclareFieldAlias{eprint:mitp}{eprint:MITP}
\hep@indico{LIP}{indico.lip.pt}
\DeclareFieldAlias{eprint:lip}{eprint:LIP}
%    \end{macrocode}
%\end{macro}
%
%</package>

% \section{Readme}
%
%<*readme>
%
%    \begin{macrocode}
# The `hep-cv-private` package

Private CV bibliography filters and environments

## Introduction

The `hep-cv-private` package adds private CV bibliography filters and
environments on top of the HEP CV bibliography setup. It groups work,
visits, education, supervision, funding, talks, publications, software,
theses, and service entries by keyword, provides specialised CV
bibliography environments for those groups, and formats Indico event and
contribution links as eprint data.

It can be loaded with `\usepackage{hep-cv-private}`.

## Author

Jan Hajer

## License

This file may be distributed and/or modified under the conditions of the
`LaTeX` Project Public License, either version 1.3c of this license or
(at your option) any later version. The latest version of this license is
in `http://www.latex-project.org/lppl.txt` and version 1.3c or later is
part of all distributions of LaTeX version 2005/12/01 or later.
%    \end{macrocode}
%
%</readme>
%
% \Finale
%
\endinput
